% Stage 8 section framing — 2.7 Non-Radiating Sources and Null Spaces
\section{Non-Radiating Sources and Null Spaces}
\label{sec:ch27}
Spectral theory in Section~\ref{sec:ch26} ranked export directions and recoverable rank.
Section~\ref{sec:ch27} addresses the complementary laboratory question: \emph{How can current
exist without observable radiation?} \cite{devaney2004}




Laboratory signatures. Three configurations recur before any kernel is displayed. A balanced wire loop in vacuum,
driven at fixed \(\omega\) with vanishing net electric dipole moment, stores reactive \(\E\) and
\(\H\) within \(\lambda_0\) of the conductor while a qualified wattmeter on \(\Gamma_R\) at
\(R\gg\lambda_0\) records zero outgoing transverse flux once \(\Pi_{\mathrm{rad}}\) selection is
enforced \cite{stratton1941}. A counter-directed dipole pair separated by
\(d\ll\lambda_0\) can annihilate leading \(1/r\) transverse fields on \(S^{2}\) while reactive
overlap remains in the gap \cite{jackson1999}. A waveguide section below cutoff carries standing
reactive fields at ports yet contributes no propagating sheet to exterior flux meters
\cite{stratton1941}. In each case Maxwell permits substantial local structure while
\(\mathcal{R}_\omega[\Jimp]=\mathbf{0}\): rich near-zone data, null far-zone export---not a
solver artifact \cite{jackson1999}.
The \emph{radiation kernel} collects every admissible impressed current whose outgoing
radiative image vanishes: currents that reorganize near-zone energy without qualified far-zone
flux \cite{devaney2004}. Near-field scanners sample
\(\mathcal{S}_\omega[\Jimp]\) on enclosing shells; far-zone flux meters and angular integrators
read only after \(\Pi_{\mathrm{rad}}\) has removed non-export content
\cite{hansen1988}. A source invisible to the latter may remain vivid to the
former---the central diagnostic that separates reactive storage from qualified export. Campaigns
that report singular values without kernel awareness merge synthesis with invisible modes
\cite{jackson1999}.
By the \emph{radiation null space} \(\ker\mathcal{R}_\omega\) we mean precisely this kernel:
admissible impressed currents whose radiative image vanishes under the outgoing map
\(\mathcal{R}_\omega=\Pi_{\mathrm{rad}}\circ\mathcal{S}_\omega\)---physically realizable
sources that store energy locally yet export no qualified far-zone power
\cite{devaney2004}. Membership in \(\ker\mathcal{R}_\omega\) is an export
statement; it does not require \(\mathcal{S}_\omega[\Jimp]=\mathbf{0}\). The balanced loop
above is a constructive instance: \(\mathcal{R}_\omega[\Jimp^{\mathrm{(loop)}}]=\mathbf{0}\)
while \(\mathcal{S}_\omega[\Jimp^{\mathrm{(loop)}}]\) remains measurable on near probes
\cite{devaney2004}.
Non-radiating sources are the central electromagnetic proof that existence does not imply
observability. A current distribution may generate substantial near-zone $\E$ and $\H$ while
exporting zero far-zone power: $\mathcal{R}_\omega$ is not injective, and
$\ker\mathcal{R}_\omega$ is generally infinite-dimensional on open domains. Different impressed
currents may therefore produce identical radiative images, angular spectra, and reciprocity
readouts while differing in reactive storage and evanescent overlap. Hidden electromagnetic
structure is not a modelling defect but a consequence of Maxwell propagation from compact
sources through outgoing selection to finite instruments \cite{stratton1941}. The absence of
radiated signatures does not imply the absence of electromagnetic structure; it identifies
source configurations whose field content remains inaccessible to far-zone observation.
Section~\ref{sec:ch26} closed the spectral chain
Eq.~\eqref{eq:ch2.section26-spectral-chain}: radiation eigenmodes, rotational degeneracy,
Fredholm solvability, and geometric filters on admissible modal sectors
\cite{stratton1941}. That program organized
\(\mathcal{R}_\omega[\Jimp]\) in modal coordinates and explained which far-zone patterns are
reachable, degenerate, or symmetry-forbidden. Section~\ref{sec:ch27} addresses the
complementary \emph{source-side} question named above: which impressed currents produce no radiative
export at all, how those currents relate to evanescent and non-propagating field content,
and what null-space structure implies for inverse problems and angular spectra
\cite{stratton1941}. Eq.~\eqref{eq:ch2.ker-R-omega} from
Subsection~\ref{sec:ch244} named \(\ker\mathcal{R}_\omega\); the subsections below develop
constructive catalog, canonical non-radiating configurations, and reconstructability
limits without re-opening Green construction, Sommerfeld completion, or radiation integrals from
Sections~\ref{sec:ch23}--\ref{sec:ch25} \cite{stratton1941}. Reciprocity constraints on
coupled modes belong to Section~\ref{sec:ch28}; plane-wave angular-spectrum action to
Section~\ref{sec:ch29}; explicit harmonic bases to \ref{ch:03}
\cite{harrington2001}. Nullity, $\mathcal{E}_{\mathrm{acc}}^{(N)}$, and realizability. Three deficits are often merged in inverse problems and structured-beam reporting
\cite{devaney2004}. \emph{Nullity} is source-side membership in
\(\ker\mathcal{R}_\omega\): far data cannot detect additions satisfying
Eq.~\eqref{eq:ch2.angular-spectrum-nr-invariance} \cite{reed1950}.
\emph{probe rank} is field-side: which directions in the angular spectrum survive
geometric, transport, and regime qualifications even when \(\mathcal{R}_\omega[\Jimp]\neq\mathbf{0}\)
\cite{harrington2001}. \emph{$\mathcal{S}_{\mathrm{real}}$}---whether a target pattern is
\(\mathcal{R}_\omega[\Jimp]\) for some admissible \(\Jimp\)---is not decided by rank-$N$ recovery
alone; \ref{ch:04} develops synthesizable proofs while Section~\ref{sec:ch27} supplies the
forward null-space catalog that constrains every claim \cite{colton1992}.
A source may store reactive energy nearby yet export no far-zone power---the radiation map annihilates whole families of impressed currents. Field-level
uniqueness after outgoing radiation condition (exterior radiative uniqueness (Subsection~\ref{sec:ch243})) does
not imply source identifiability: many admissible \(\Jimp\in\mathcal{J}_\Omega\) can share
the same \(\mathcal{R}_\omega[\Jimp]\) through Eq.~\eqref{eq:ch2.ker-R-omega}
\cite{stratton1941}. Spectral mode \(\lambda_n=0\) in
Eq.~\eqref{eq:ch2.modal-eigenvalue-equation} marks a radiative eigen-direction with
vanishing export, yet \(\ker\mathcal{R}_\omega\) is generally infinite-dimensional on open
domains and is not exhausted by finitely many modal labels
\cite{harrington2001}. Using cavity modal books or truncated harmonic sums to
\emph{approximate} \(\ker\mathcal{R}_\omega\) misstates identifiability even when outgoing
source-to-field integrals from Subsection~\ref{sec:ch251} are evaluated exactly
\cite{jackson1999}. Section~\ref{sec:ch27} separates operator null spaces,
field-side non-radiative components Eqs.~\eqref{eq:ch2.radiative-component-def}--\eqref{eq:ch2.nonradiative-component-def},
and observation null spaces Eq.~\eqref{eq:ch2.observation-null-space}, then forwards
reciprocity and angular-spectrum filters to Sections~\ref{sec:ch28}--\ref{sec:ch29}
\cite{stratton1941}. Let \(\Omega_s=\operatorname{supp}\Jimp\) be compact, \(\omega>0\) fixed, and outgoing field
pairs lie in \(\mathcal{F}_{\mathrm{SF}}\) from
Eq.~\eqref{eq:ch2.admissible-field-space}. Section~\ref{sec:ch27} organizes analysis in four
stages:
\begin{equation}
\label{eq:ch2.section27-nullspace-chain}
\ker\mathcal{R}_\omega
\ \xrightarrow{\ \text{operator structure}\ }
\text{null-space bases}
\ \xrightarrow{\ \text{canonical sources}\ }
\text{NR current families}
\ \xrightarrow{\ \text{evanescent overlap}\ }
\text{non-reconstructability}
\ \xrightarrow{\ \text{angular spectra}\ }
\text{suppressed sheets},
\end{equation}
where the first arrow is Subsection~\ref{sec:ch271}, the second is Subsection~\ref{sec:ch272},
the third is Subsection~\ref{sec:ch273}, and the fourth is Subsection~\ref{sec:ch274}
\cite{stratton1941}. Equation~\eqref{eq:ch2.section27-nullspace-chain}
expresses Eq.~\eqref{eq:ch2.section26-spectral-chain} at the source level: spectral
\(\lambda_n=0\) modes from Eq.~\eqref{eq:ch2.spectral-resolution-composition} lie in
\(\mathbf{u}_{\mathrm{nr}}\), but complete \(\ker\mathcal{R}_\omega\) requires impressed-data
coordinates on \(\mathcal{J}_\Omega\), not only far-zone modal shapes
\cite{stratton1941}.
Equation~\eqref{eq:ch2.section27-nullspace-chain} is the source-side restatement
of the spectral chain Eq.~\eqref{eq:ch2.section26-spectral-chain}: operator
nullity precedes constructive NR catalogs, which precede evanescent
non-reconstructability, which precede angular-spectrum consequences.
Mathematically, each arrow preserves linearity of \(\mathcal{R}_\omega\);
physically, each stage removes a distinct laboratory error from inverse-source
charting. On a cube of side \(2a\) with \(k_0 a=2\), evanescent overlap with
\(\Pi_{\mathrm{ev}}\mathcal{S}_\omega[\Jimp]\) can exceed \(25\%\) of the
near-field energy on \(S_R\) with \(R=a\) while contributing zero to
\(a(\widehat{\mathbf{k}})\); mistaking that storage for NR nullity violates
Eq.~\eqref{eq:ch2.evanescent-nonexport-law} \cite{stratton1941}.
\ref{ch:01} flux laws enter only after \(\Pi_{\mathrm{rad}}\)
\cite{harrington2001}.
Composition law Eq.~\eqref{eq:ch2.constraint-composition-law} still
applies: geometric annihilation Eq.~\eqref{eq:ch2.geometric-symmetry-operator-def} and
observation rank Eq.~\eqref{eq:ch2.reconstruction-rank-cap} are independent of, and
composable with, source null-space membership \cite{harrington2001}.
The laboratory examples and kernel reading above fix the measurement ordering before any
operator display. Near-field probes sample \(\mathcal{S}_\omega[\Jimp]\); flux meters read only
after \(\Pi_{\mathrm{rad}}\) \cite{stratton1941}. With that ordering in place,
keep three spaces in view throughout Section~\ref{sec:ch27}: impressed currents
\(\Jimp\in\mathcal{J}_\Omega\) on compact support \(\Omega_s\), full Maxwell fields
\(\mathcal{S}_\omega[\Jimp]\in\mathcal{F}_\Omega\), and the radiative image on which far-zone
instruments read. The primary object remains the radiation map that projects the latter from the
former:
\begin{equation}
\label{eq:ch2.nullspace-operator-targets}
\mathcal{R}_\omega:\mathcal{J}_\Omega\longrightarrow\mathcal{F}_{\mathrm{rad}},
\qquad
\mathcal{R}_\omega[\Jimp]=\Pi_{\mathrm{rad}}\,\mathcal{S}_\omega[\Jimp],
\end{equation}
from Eqs.~\eqref{eq:ch2.radiation-operator-def}--\eqref{eq:ch2.outgoing-radiation-refinement}
\cite{stratton1941}. The displayed map is the mathematical encoding of the wattmeter reading:
only \(\Pi_{\mathrm{rad}}\mathcal{S}_\omega[\Jimp]\) enters \(\mathcal{F}_{\mathrm{rad}}\), so
near-zone energy in \(\mathcal{S}_\omega[\Jimp]\setminus\mathcal{F}_{\mathrm{rad}}\) cannot be
inferred from far data alone \cite{stratton1941}. Secondary targets include the observation map
\(\mathcal{M}_\omega=\mathcal{O}_{\mathrm{meas}}\circ\Lambda_{\mathrm{rad}}\) from
Eq.~\eqref{eq:ch2.reconstruction-observation-map}, the non-radiative field projector
\(\mathrm{id}-\Pi_{\mathrm{rad}}\) on \(\mathcal{F}_\Omega\), and evanescent/non-propagating
selectors Eqs.~\eqref{eq:ch2.spectral-projectors}--\eqref{eq:ch2.nonpropagating-component-def}
from Subsection~\ref{sec:ch224} \cite{stratton1941}. Inverse source Eq.~\eqref{eq:ch2.inverse-source-equation}
is ill-posed whenever \(\ker\mathcal{M}_\omega\neq\{\mathbf{0}\}\), which occurs if
\(\ker\mathcal{R}_\omega\neq\{\mathbf{0}\}\) or \(\ker\mathcal{O}_{\mathrm{meas}}\) is
nontrivial on \(\operatorname{ran}\mathcal{R}_\omega\) \cite{harrington2001}.
reciprocity-based solvability Eq.~\eqref{eq:ch2.radiation-solvability-criterion} governs reachability
of a \emph{target} radiative field; Section~\ref{sec:ch27} classifies sources that map to
\(\mathbf{0}\in\mathcal{F}_{\mathrm{rad}}\) regardless of solvability for nonzero targets
\cite{stratton1941}.
Radiation null space (leading template). Subsection~\ref{sec:ch271} develops Eq.~\eqref{eq:ch2.ker-R-omega} to
structural statements on \(\ker\mathcal{R}_\omega\) and \(\ker\mathcal{R}_\omega^{\dagger}\).
Locally, \(\ker\mathcal{R}_\omega\) collects impressed currents annihilated by export;
\(\ker\mathcal{R}_\omega^{\dagger}\) collects far-side test directions orthogonal to every
radiative image; and \(\mathrm{coker}\,\mathcal{R}_\omega=\mathcal{F}_{\mathrm{rad}}/\operatorname{ran}\mathcal{R}_\omega\)
labels pattern targets that no admissible source can realize \cite{stratton1941}. The leading template
below records admissibility and typical infinite dimensionality on open domains:
\begin{equation}
\label{eq:ch2.radiation-nullspace-structure}
\ker\mathcal{R}_\omega
=
\left\{
\Jimp\in\mathcal{J}_\Omega:
\Pi_{\mathrm{rad}}\,\mathcal{S}_\omega[\Jimp]=\mathbf{0},
\ \mathcal{D}_{J}[\Jimp]=1
\right\},
\qquad
\dim\ker\mathcal{R}_\omega=\infty\ \text{typ. on open domains},
\end{equation}
with $\mathcal{D}_J=1$ Eq.~\eqref{eq:ch2.source-admissibility} enforced throughout
\cite{stratton1941}. Equation~\eqref{eq:ch2.radiation-nullspace-structure}
does not re-define Eq.~\eqref{eq:ch2.ker-R-omega}; it records the Chapter~\ref{ch:02}
primary target for constructive characterization of nullity, dimensionality, and pairing with cokernel
directions obstructing inversion Eq.~\eqref{eq:ch2.spectral-solvability}
\cite{stratton1941}. Modal sectors with \(\lambda_n=0\) in
Eq.~\eqref{eq:ch2.spectral-resolution-composition} supply a finite-dimensional \emph{slice}
of \(\ker\mathcal{R}_\omega\) only when eigenmodes are declared on outgoing subspaces
\cite{harrington2001}.
Non-radiating current configurations (leading template). Subsection~\ref{sec:ch272} develops canonical impressed currents that realize
Eq.~\eqref{eq:ch2.null-radiation}:
\begin{equation}
\label{eq:ch2.nr-current-template}
\Jimp^{\mathrm{(NR)}}\in\ker\mathcal{R}_\omega,
\qquad
\mathcal{S}_\omega[\Jimp^{\mathrm{(NR)}}]\neq\mathbf{0},
\qquad
\|\mathbf{u}_{\mathrm{nr}}\|\ \text{may be large near}\ \Omega_s,
\end{equation}
with \(\mathbf{u}_{\mathrm{nr}}\) from Eq.~\eqref{eq:ch2.nonradiative-component-def}
\cite{stratton1941}. Equation~\eqref{eq:ch2.nr-current-template}
specializes Eq.~\eqref{eq:ch2.nr-image-of-null-space} to constructive families---loop modes,
counter-propagating pairs, and equivalent-surface cancellations compatible with
Eq.~\eqref{eq:ch1.charge-continuity-harmonic}---without tabulating dyadic kernels from
Subsection~\ref{sec:ch231} \cite{stratton1941}. Such currents are invisible to far-zone flux
meters Eq.~\eqref{eq:ch2.power-flux-functional} yet may dominate near-zone energy budgets
Eq.~\eqref{eq:ch2.component-power-budget} \cite{harrington2001}.
Evanescent overlap and non-reconstructability (leading template). Subsection~\ref{sec:ch273} links \(\ker\mathcal{R}_\omega\) to evanescent subspaces and
finite observation rank. Evanescent content is field-side reactive storage with
\(\Pi_{\mathrm{rad}}\mathbf{u}_{\mathrm{ev}}=\mathbf{0}\); it must not be conflated with
source-side \(\Jimp\in\ker\mathcal{R}_\omega\), although both deficits can coexist on the same
excitation \cite{stratton1941}. The relation records the measurement inclusion and
the evanescent non-export law:
\begin{equation}
\label{eq:ch2.evanescent-nonreconstruct}
\ker\mathcal{O}_{\mathrm{meas}}\circ\mathcal{R}_\omega
\supseteq
\ker\mathcal{R}_\omega,
\qquad
\Pi_{\mathrm{ev}}\mathcal{S}_\omega[\Jimp]
\ \text{never exports through}\ \Pi_{\mathrm{rad}},
\end{equation}
using Eqs.~\eqref{eq:ch2.rad-null-on-evanescent} and~\eqref{eq:ch2.reconstruction-rank-cap}
\cite{stratton1941}. Equation~\eqref{eq:ch2.evanescent-nonreconstruct}
separates \emph{radiation nullity} (no outgoing transverse export) from \emph{measurement
nullity} (unobservable even when export is nonzero) \cite{stratton1941}. Tikhonov stabilization
Eq.~\eqref{eq:ch2.tikhonov-source-reconstruction} does not shrink \(\ker\mathcal{R}_\omega\);
regularization only damps ill-conditioned directions on \(\operatorname{ran}\mathcal{R}_\omega\)
\cite{harrington2001}. Accessible modal sector
Eq.~\eqref{eq:ch2.accessible-modal-sector} from Subsection~\ref{sec:ch264} further restricts
recoverable content independently of evanescent storage \cite{stratton1941}.
Angular-spectrum consequences (leading template). Subsection~\ref{sec:ch274} records how null-space structure propagates to continuous angular
spectra Eq.~\eqref{eq:ch2.exterior-continuous-spectral-integral}:
\begin{equation}
\label{eq:ch2.angular-spectrum-null}
a(\widehat{\mathbf{k}})=0
\ \text{on suppressed sheets}\ \mathcal{K}_{\mathrm{sup}}
\ \Longrightarrow\
\Jimp\ \text{may still carry reactive content in}\ \Pi_{\mathrm{ev}}\mathcal{S}_\omega[\Jimp],
\end{equation}
with \(a(\widehat{\mathbf{k}})\) the far angular amplitude and
\(\mathcal{K}_{\mathrm{sup}}\) the complement of accessible sector
Eq.~\eqref{eq:ch2.angular-accessibility-sector} \cite{stratton1941}. Equation
\eqref{eq:ch2.angular-spectrum-null} links Section~\ref{sec:ch27} to
Section~\ref{sec:ch29}: plane-wave filters and retained/suppressed components in
Subsections~\ref{sec:ch291}--\ref{sec:ch294} presuppose the
source-side null-space catalog developed here \cite{stratton1941}. Angular probe rank on the far shell (Subsection~\ref{sec:ch264}) is not re-proved here;
Subsection~\ref{sec:ch274} states angular-spectrum \emph{consequences} of
Eq.~\eqref{eq:ch2.ker-R-omega} \cite{harrington2001}.
\begin{equation}
\label{eq:ch2.nullspace-spectral-composition}
\Jimp\in\ker\mathcal{R}_\omega
\ \Longrightarrow\
a_n(\Jimp)=0\ \text{for all exporting modes with}\ \lambda_n\neq 0,
\end{equation}
while geometric constraints Eq.~\eqref{eq:ch2.accessible-modal-sector} may force additional
\(a_n=0\) for \(\Jimp\notin\ker\mathcal{R}_\omega\) \cite{stratton1941}.
Equation~\eqref{eq:ch2.nullspace-spectral-composition} follows from
Eq.~\eqref{eq:ch2.spectral-far-pattern-expansion} and linearity
Eq.~\eqref{eq:ch2.radiation-linearity}; it is cited in Section~\ref{sec:ch27}
\cite{stratton1941}. Inverse problems must budget \(\ker\mathcal{R}_\omega\),
\(\ker\mathcal{O}_{\mathrm{meas}}\), and forbidden sectors from
geometric constraints on radiative spectra (Subsection~\ref{sec:ch264}) separately
\cite{harrington2001}.
Several constructions deliberately remain outside Section~\ref{sec:ch27}. Reciprocity
identities and modal orthogonality belong to Section~\ref{sec:ch28}; plane-wave angular-spectrum
operators and near-to-far transforms to Section~\ref{sec:ch29}; vector spherical harmonic
completeness and explicit \(\mathbf{f}_n\) to \ref{ch:03}; range membership of
unreachable modes to \ref{ch:04} \cite{stratton1941}. Section~\ref{sec:ch27}
Subsection~\ref{sec:ch261} supplies spectral decomposition;
Fredholm structure of radiation operators (Subsection~\ref{sec:ch263}), or geometric constraints
on radiative spectra (Subsection~\ref{sec:ch264}); it classifies invisible sources on the
maps those results presuppose \cite{stratton1941}. Worked non-radiating examples appear
in Subsection~\ref{sec:ch2102} by reference to Subsection~\ref{sec:ch272}
\cite{harrington2001}.
Regime logic from Section~\ref{sec:ch01.2.4} is unchanged. Null-space membership does not
override transport qualification: a current in \(\ker\mathcal{R}_\omega\) has
\(\mathcal{P}_{\mathrm{rad}}=0\) on qualified shells, while reactive near fields from
\(\mathcal{S}_\omega[\Jimp]\) may still violate naive flux interpretations if
Eq.~\eqref{eq:ch2.regime-operator-screen} is ignored \cite{stratton1941}.
Far-zone dominance far-zone radiative dominance (Subsection~\ref{sec:ch254}) applies only to
exporting sources; it does not trivialize \(\ker\mathcal{R}_\omega\) \cite{stratton1941}.
Compression error Eq.~\eqref{eq:ch2.compression-error-bound} bounds modal truncation on
\(\operatorname{ran}\mathcal{R}_\omega\) independently of null-space dimension
\cite{harrington2001}.
As a worked example, consider a thin rectangular wire loop of perimeter \(2\pi b\) with \(b\ll\lambda_0\), embedded in
free space and driven at angular frequency \(\omega\) so that the impressed current completes one
closed circulation with equal and opposite arms along the loop's long sides. The physical setup is
a balanced dipole pair folded into a single closed contour: magnetic flux linkage is negligible at
\(\omega\), yet the loop stores reactive energy in the near zone. A closed admissible loop carrying
oppositely directed arms can cancel far-zone transverse radiation through
Eq.~\eqref{eq:ch2.radiation-integral-outgoing-upgrade} while maintaining nonzero
\(\mathcal{S}_\omega[\Jimp]\) near the loop \cite{stratton1941}.
Subsection~\ref{sec:ch272} classifies such families; here the point is that
Eq.~\eqref{eq:ch2.null-radiation} is realizable with
\(\mathcal{D}_{J}[\Jimp]=1\), not merely a formal kernel direction
\cite{harrington2001}. Pattern-matching algorithms that fit only
\(\operatorname{ran}\mathcal{R}_\omega\) assign spurious modal labels to invisible loop content
\cite{stratton1941}.
As a worked example, take a compact spherical source of radius \(a\) with \(k_0 a=1.5\) and truncate the exterior modal
book at \(N=12\) radiation eigenpairs. A single \(\lambda_n=0\) eigenmode from
Eq.~\eqref{eq:ch2.radiation-eigenpair-def} contributes to \(\ker\mathcal{R}_\omega\) when
realized by admissible \(\Jimp^{(n)}\) \cite{stratton1941}. On the range, the truncated
index \(n\le N\) labels only a finite slice of invisible degrees; the full kernel on open
\(\mathcal{J}_\Omega\) is generally much larger. On exterior problems the modal index set is not finite;
truncating to the first \(N\) modes underestimates \(\dim\ker\mathcal{R}_\omega\) and
overstates identifiability of inverse Eq.~\eqref{eq:ch2.inverse-source-equation}
\cite{harrington2001}. Subsection~\ref{sec:ch271} states the correct operator
bookkeeping \cite{stratton1941}.
